% =============================================================================
% Ringkasan Visual dan Graphical Summary -- satu halaman
% Gambar akhir disarankan memiliki resolusi sekurang-kurangnya 300 dpi.
% =============================================================================

\ipbprelimheading{RINGKASAN VISUAL}

\begin{center}
\ifdefempty{\BerkasRingkasanVisual}{%
  \fbox{\parbox[c][5.7cm][c]{0.8\textwidth}{\centering
    Sisipkan sebuah gambar yang secara visual merangkum temuan utama}}%
}{%
  \includegraphics[width=0.8\textwidth,height=5.7cm,keepaspectratio]{\BerkasRingkasanVisual}%
}

\vspace{0.3cm}

{\bfseries\parbox{9cm}{\centering\JudulRingkasanVisual}\par}
\end{center}

\smallskip

\begingroup
\setlength{\parindent}{1.27cm}%
Ringkasan visual/\textit{graphical summary} terdiri atas sebuah gambar yang
secara visual merangkum temuan utama dari tugas akhir. Ringkasan visual ini
disusun secara singkat dan berdiri sendiri, menyampaikan tujuan dan hasil
penelitian, serta disajikan dengan cara yang menarik secara visual dengan
resolusi gambar minimal 300 dpi. Ringkasan visual/\textit{graphical summary}
ditulis dalam satu halaman. Judul diletakkan di bawah ringkasan
visual/\textit{graphical summary}. Mahasiswa wajib mendeklarasikan penggunaan
AI jika digunakan dalam pembuatan ringkasan visual dan perhatikan penggunaan
elemen berhak cipta (lihat Suplemen 6).
\par\endgroup

\begin{otherlanguage}{english}
\itshape
\ipbprelimsubheading{GRAPHICAL SUMMARY}

\begin{center}

\ifdefempty{\BerkasRingkasanVisualInggris}{%
  \fbox{\parbox[c][5.7cm][c]{0.8\textwidth}{\centering
    Insert a single image that visually summarizes the main findings}}%
}{%
  \includegraphics[width=0.8\textwidth,height=5.7cm,keepaspectratio]{\BerkasRingkasanVisualInggris}%
}

\vspace{0.3cm}

{\bfseries\parbox{9cm}{\centering\JudulRingkasanVisualInggris}\par}
\end{center}
\end{otherlanguage}
